\subsection{Marginal diagnostics and data-quality strata}
Figure~\ref{fig:bias} shows that dependence comparisons inherit distinct
marginal errors. At 60 minutes, mean Borg pending-task exit probability is
0.3193 versus an observed fraction of 0.3184. For running tasks the corresponding
values are 0.1199 and 0.1287. Their final integrated Brier scores are 0.08251
and 0.02937. The BPI 14-day bias is larger in absolute fraction, as discussed
above. A small pooled mean error is not evidence of conditional calibration at
every origin or relation group; the released reliability and origin tables
retain this variation.

Cutoff-defined strata also matter. About 21,876 running tasks per final origin
have unknown initial state age. Their integrated Brier is 0.00025, compared
with 0.03529 for known-age running tasks. Long-lived initial-state entities are
thus a relatively easy subpopulation, and pooled scores conceal that
heterogeneity. The unknown-age pending stratum contains only one task per
origin and cannot support a general robustness claim. Running tasks with a
known preceding publisher-imputed event have Brier 0.01078 versus 0.03006
without one; these are descriptive case-mix differences, not causal effects
of imputation. All strata remain in the primary complete-cohort evaluation.

\subsection{Numerical sensitivity and execution}
Doubling seed-0 scenario draws from 512 to 1,024 changes macro CRPS by at
most 0.000122 across the two applications, families, and five controls.
Within each application and family, the ordering of the five controls is
unchanged. This checks the reported CRPS comparisons against finite-ensemble
variation; it does not replace uncertainty over future forecast origins.

On the shared RTX A6000 host, each 512-draw Borg evaluation of all five
controls across 28 origins takes 245.8--251.0
seconds wall time, including sampling and scoring. The corresponding BPI
evaluation across 61 origins takes 3.48--3.68
seconds. Borg peak allocated PyTorch memory is at most 0.67\,GiB
per evaluation process and peak host resident memory is 2.05\,GiB.
These are evaluation costs with individual CDFs already computed, measured
while other evaluations share the device; they exclude data reconstruction,
training, and individual-model inference. They are not an isolated latency
benchmark.

The final artifact audit verifies all 264 frozen
file hashes, every final CDF and label trajectory, the exact frozen feature
transform, and the expected origins and comparisons. No frozen model or
evaluation file was modified during final evaluation. The separately audited
lossless timestamp correction is described in the appendix.

